\begin{exercice*}
    Voici un programme.

    \begin{Scratch}[Echelle=0.9]
        drapeau
        demander("Choisis un nombre")
        mettrevariable("Résultat", add(mul(3,réponse),7))
        mettrevariable("Résultat", mul(variable(Résultat), réponse))
        mettrevariable("Résultat", sous(variable(Résultat), réponse))
        dire(variable(Résultat))
    \end{Scratch}
    \begin{enumerate}
        \item Détermine le résultat obtenu pour le nombre de ton choix.
        \item Exprime le résultat obtenu par le programme pour un nombre $x$.
        \item Charles remarque qu'en choisissant un nombre entier, le programme donne toujours un multiple de 3. Justifie cette remarque.
    \end{enumerate}
\end{exercice*}
\begin{corrige}
    %\setcounter{partie}{0} % Pour s'assurer que le compteur de \partie est à zéro dans les corrigés
    %\phantom{rrr}    
    \dots
\end{corrige}

